% Nguồn SGK Toán 9 Tập một — Bài tập cuối chương III, trang 65:
% https://cdn3.olm.vn/upload/thuvienso/4491668113-page-65-1773022942147.png
%
% Nguồn SBT Toán 9 Tập một — Ôn tập chương III, trang 39--40:
% SBT TOAN 9 KNTT TAP 1.pdf
% SHA-256: d7ffd04618456682795977e5c47de7874161a35e8606e5bac1906881603a72f7
% Không lấy phần lời giải/hướng dẫn/đáp số của SBT.

\section*{Bài tập cuối chương III}

\begin{exercises}
    \subsection{Bài tập cuối chương III}

    \textbf{A. TRẮC NGHIỆM}

    \begin{questions}[resume=SGK]
        \item Căn bậc hai của $4$ là
        \begin{tasks}(4)
            \task[(A)] $2$.
            \task[(B)] $-2$.
            \task[(C)] $2$ và $-2$.
            \task[(D)] $\sqrt{2}$ và $-\sqrt{2}$.
        \end{tasks}
        \par\noindent\answerbox\par

        \item Căn bậc hai số học của $49$ là
        \begin{tasks}(4)
            \task[(A)] $7$.
            \task[(B)] $-7$.
            \task[(C)] $7$ và $-7$.
            \task[(D)] $\sqrt{7}$ và $-\sqrt{7}$.
        \end{tasks}
        \par\noindent\answerbox\par

        \item Rút gọn biểu thức $\sqrt[3]{\left(4-\sqrt{17}\right)^3}$ ta được
        \begin{tasks}(2)
            \task[(A)] $4+\sqrt{17}$.
            \task[(B)] $4-\sqrt{17}$.
            \task[(C)] $\sqrt{17}-4$.
            \task[(D)] $-4-\sqrt{17}$.
        \end{tasks}
        \par\noindent\answerbox\par

        \item Độ dài đường kính (mét) của hình tròn có diện tích $4\,\mathrm{m}^2$ sau khi làm tròn kết quả đến chữ số thập phân thứ hai bằng
        \begin{tasks}(4)
            \task[(A)] $2{,}26$.
            \task[(B)] $2{,}50$.
            \task[(C)] $1{,}13$.
            \task[(D)] $1{,}12$.
        \end{tasks}
        \par\noindent\answerbox\par

        \item Một vật rơi tự do từ độ cao $396{,}9$ m. Biết quãng đường chuyển động $S$ (mét) của vật phụ thuộc vào thời gian $t$ (giây) bởi công thức $S=4{,}9t^2$. Vật chạm đất sau
        \begin{tasks}(4)
            \task[(A)] $8$ giây.
            \task[(B)] $5$ giây.
            \task[(C)] $11$ giây.
            \task[(D)] $9$ giây.
        \end{tasks}
        \par\noindent\answerbox\par
    \end{questions}

    \textbf{B. TỰ LUẬN}

    \begin{questions}[resume=SGK]
        \item Không sử dụng MTCT, tính giá trị của biểu thức
        \[
            A=\sqrt{\left(\sqrt{3}-2\right)^2}
            +\sqrt{4\left(2+\sqrt{3}\right)^2}
            -\dfrac{1}{2-\sqrt{3}}.
        \]
        \answerline
        \answerline

        \item Cho biểu thức
        \[
            A=\dfrac{\sqrt{x}+2}{\sqrt{x}-2}
            -\dfrac{4}{\sqrt{x}+2}\quad (x\ge0\ \text{và}\ x\ne4).
        \]
        \begin{questions}
            \item Rút gọn biểu thức $A$.
            \item Tính giá trị của $A$ tại $x=14$.
        \end{questions}
        \answerline
        \answerline

        \item Biết rằng nhiệt lượng tỏa ra trên dây dẫn được tính bởi công thức $Q=I^2Rt$, trong đó $Q$ là nhiệt lượng tính bằng đơn vị Joule (J), $R$ là điện trở tính bằng đơn vị Ohm ($\Omega$), $I$ là cường độ dòng điện tính bằng đơn vị Ampe (A), $t$ là thời gian tính bằng giây (s). Dòng điện chạy qua một dây dẫn có $R=10\,\Omega$ trong thời gian $5$ giây.
        \begin{questions}
            \item Thay dấu ``?'' trong bảng sau bằng các giá trị thích hợp.

            \begin{center}
                \begin{tabular}{|c|c|c|c|}
                    \hline
                    $I$ (A) & $1$ & $1{,}5$ & $2$ \\
                    \hline
                    $Q$ (J) & ? & ? & ? \\
                    \hline
                \end{tabular}
            \end{center}

            \item Cường độ dòng điện là bao nhiêu Ampe để nhiệt lượng tỏa ra trên dây dẫn đạt $800$ J?
        \end{questions}
        \answerline
        \answerline
    \end{questions}
\end{exercises}

\section*{Ôn tập chương III}

\begin{practice}
    \subsection{Ôn tập chương III}

    \textbf{A. CÂU HỎI (Trắc nghiệm)}

    \begin{enumerate}
        \item Trong các khẳng định sau, khẳng định nào đúng?
        \begin{tasks}(1)
            \task[(A)] Mọi số thực đều có căn bậc hai.
            \task[(B)] Các số thực âm không có căn bậc ba.
            \task[(C)] Mọi số thực đều có căn bậc hai số học.
            \task[(D)] Các số thực dương có hai căn bậc hai đối nhau.
        \end{tasks}
        \par\noindent\answerbox\par

        \item Trong các khẳng định sau, khẳng định nào sai?
        \begin{tasks}(1)
            \task[(A)] Mọi số thực dương đều có căn bậc hai.
            \task[(B)] Số $-216$ không có căn bậc ba.
            \task[(C)] Số $-216$ không có căn bậc hai.
            \task[(D)] Số $-216$ không có căn bậc hai số học.
        \end{tasks}
        \par\noindent\answerbox\par

        \item Cho biết $3{,}1^2=9{,}61$. Số nào sau đây là giá trị của $\sqrt{0{,}000961}$?
        \begin{tasks}(4)
            \task[(A)] $3{,}1$.
            \task[(B)] $0{,}31$.
            \task[(C)] $0{,}031$.
            \task[(D)] $0{,}000031$.
        \end{tasks}
        \par\noindent\answerbox\par

        \item Biểu thức nào sau đây có giá trị khác với các biểu thức còn lại?
        \begin{tasks}(4)
            \task[(A)] $\sqrt{7^2}$.
            \task[(B)] $\sqrt{(-7)^2}$.
            \task[(C)] $(-\sqrt{7})^2$.
            \task[(D)] $-(\sqrt{7})^2$.
        \end{tasks}
        \par\noindent\answerbox\par

        \item Độ dài cạnh khối lập phương có thể tích bằng $0{,}512\,\mathrm{dm}^3$ là
        \begin{tasks}(4)
            \task[(A)] $8$ cm.
            \task[(B)] $8$ dm.
            \task[(C)] $0{,}8$ cm.
            \task[(D)] $0{,}08$ dm.
        \end{tasks}
        \par\noindent\answerbox\par
    \end{enumerate}

    \textbf{B. BÀI TẬP}

    \begin{questions}[resume=SBT]
        \item Thực hiện phép tính:
        \begin{questions}
            \item $\sqrt{12-\sqrt{23}}\cdot\sqrt{12+\sqrt{23}}$;
            \item $\left(\sqrt{9-\sqrt{17}}+\sqrt{9+\sqrt{17}}\right)^2$.
        \end{questions}
        \answerline
        \answerline

        \item So sánh $\sqrt{\sqrt{89+24\sqrt{5}}}$ và $\sqrt{1+\sqrt{122}}$.
        \answerline
        \answerline

        \item
        \begin{questions}
            \item Chứng minh rằng
            \[
                \sqrt{3+\sqrt{5}}\cdot\sqrt{3-\sqrt{5}}=2
                \quad\text{và}\quad
                \sqrt{3+\sqrt{5}}+\sqrt{3-\sqrt{5}}=\sqrt{10}.
            \]
            \item Rút gọn các biểu thức sau:
            \[
                A=\left(\sqrt{3+\sqrt{5}}\right)^3
                +\left(\sqrt{3-\sqrt{5}}\right)^3;
            \]
            \[
                B=\left(\sqrt{3+\sqrt{5}}\right)^5
                +\left(\sqrt{3-\sqrt{5}}\right)^5.
            \]
        \end{questions}
        \answerline
        \answerline

        \item Cho $x$, $y$ là hai số dương thoả mãn $x^2+y^2=1$. Tính giá trị biểu thức
        \[
            A=x-y+\sqrt{1-x^2}-\sqrt{1-y^2}.
        \]
        \answerline
        \answerline

        \item
        \begin{questions}
            \item Khai triển $\left(2-\sqrt{3}\right)^2$ và $\left(2\sqrt{3}-3\right)^2$ thành những biểu thức không còn bình phương.
            \item Sử dụng kết quả câu a, rút gọn các biểu thức sau:
            \[
                A=\sqrt{4-2\sqrt{3}-\sqrt{21-12\sqrt{3}}};
            \]
            \[
                B=\sqrt{2+\sqrt{3}+\sqrt{4-2\sqrt{3}-\sqrt{21-12\sqrt{3}}}}.
            \]
        \end{questions}
        \answerline
        \answerline
    \end{questions}
\end{practice}
